% Generated by scripts/assemble.py; corrections belong in source/corrections.json.
\examhead{136}{武汉理工大学 2025 年数学分析试卷}{130}
% p131-b01; source PDF page 131
\par\Needspace{5\baselineskip}
\blockanchor{p131-b01}
\majorquestion{一}计算题.\par
% p131-b02; source PDF page 131
\blockanchor{p131-b02}
\question{1}计算极限 $\lim_{x\to 0}\frac{\sqrt[3]{1 + \tan^2x} - \sqrt{\cos x}}{x\sin x}$.\par
% p131-b03; source PDF page 131
\par\Needspace{6\baselineskip}
\blockanchor{p131-b03}
\question{2}设 $f(u, v)$ 在 $\mathbb{R}^2$ 上有二阶连续偏导数, 且 $\frac{\partial^2 f}{\partial u^2} + \frac{\partial^2 f}{\partial v^2} = 1$, 设 $g(x, y) = f\left(xy, \frac{x^2 - y^2}{2}\right)$, 求\par
\nopagebreak[4]
% p131-b04; source PDF page 131
\blockanchor{p131-b04}
\displayformula{\frac {\partial^ {2} g}{\partial x ^ {2}} + \frac {\partial^ {2} g}{\partial y ^ {2}}.}
% p131-b05; source PDF page 131
\blockanchor{p131-b05}
\question{3}计算 $\iint_{D} \frac{x \sin^2 y}{y^2} \, \mathrm{d}x \, \mathrm{d}y$, 其中 $D$ 是由 $y = x, x = 0, y = \pi$ 围成的区域.\par
% p131-b06; source PDF page 131
\blockanchor{p131-b06}
\question{4}计算 $\int_{L} y e^{y^2} \, \mathrm{d}y + (x e^{y^2} + 2xy^2 e^{y^2}) \, \mathrm{d}x$, 其中 $L: y = \sqrt[3]{x}$ 从 $(0,0)$ 到 $(1,1)$.\par
% p131-b07; source PDF page 131
\blockanchor{p131-b07}
\question{5}设曲线 $C: \left\{ \begin{array}{l} x^2 + y^2 = 2az; \\ x^2 + y^2 + xy = a^2. \end{array} \right.$ 求曲线上点到坐标系平面 $xOy$ 的最大值与最小值.\par
% p131-b08; source PDF page 131
\blockanchor{p131-b08}
\question{6}讨论级数 $\sum_{n=1}^{\infty} \frac{2^n}{n} x^{2n+1}$ 的收敛域并求其和函数.\par
% p131-b09; source PDF page 131
\par\Needspace{5\baselineskip}
\blockanchor{p131-b09}
\majorquestion{二}简答题.\par
% p131-b10; source PDF page 131
\blockanchor{p131-b10}
\question{7}设  $\varphi(x,y)$  在  $(0,0)$  的某邻域内连续， $f(x,y)=|x-y|\varphi(x,y)$ .\par
% p131-b11; source PDF page 131
\blockanchor{p131-b11}
(1) 讨论 $\varphi(x, y)$ 满足什么条件时, $f_x(0, 0), f_y(0, 0)$ 存在.\par
% p131-b12; source PDF page 131
\blockanchor{p131-b12}
(2) 讨论 $\varphi(x, y)$ 满足什么条件时, $f(x, y)$ 在 $(0, 0)$ 处可微.\par
% p131-b13; source PDF page 131
\blockanchor{p131-b13}
\question{8}讨论反常积分 $\int_{1}^{+\infty} \frac{\sin x}{x^p} \, \mathrm{d}x (p > 0)$ 的敛散性, 并说明收敛时是绝对收敛还是条件收敛.\par
% p131-b14; source PDF page 131
\par\Needspace{5\baselineskip}
\blockanchor{p131-b14}
\majorquestion{三}证明题.\par
% p131-b15; source PDF page 131
\blockanchor{p131-b15}
\question{9}设 $f(x)$ 在 $(- \infty, +\infty)$ 上一致连续, $f(x) \in [a, b]$, $g(x)$ 在 $[a, b]$ 上连续, 证明: $g(f(x))$ 在 $(- \infty, +\infty)$ 上一致连续.\par
% p131-b16; source PDF page 131
\blockanchor{p131-b16}
\question{10}证明: $f(x)$ 在 $[a, b]$ 上连续, 则 $f(x)$ 在 $[a, b]$ 上有界.\par
% p131-b17; source PDF page 131
\blockanchor{p131-b17}
\question{11}设数列 $\{a_{n}\}$ 为恒正的单调增数列, 证明:\par
% p131-b18; source PDF page 131
\blockanchor{p131-b18}
(1) 若 $\{a_{n}\}$ 有界, 则 $\sum_{n=1}^{\infty}\left(1 - \frac{a_n}{a_{n+1}}\right)$ 收敛.\par
% p131-b19; source PDF page 131
\blockanchor{p131-b19}
(2) 若 $\{a_{n}\}$ 无界, 则 $\sum_{n=1}^{\infty}\left(1 - \frac{a_n}{a_{n+1}}\right)$ 发散.\par
% p131-b20; source PDF page 131
\par\Needspace{6\baselineskip}
\blockanchor{p131-b20}
\question{12}设 $f(x)$ 在 $[0, +\infty)$ 上连续, $A > 0$, $\int_{A}^{+\infty} \frac{f(x)}{x} \mathrm{d}x$ 存在, 证明: 对任意的 $a, b > 0$, 有\par
\nopagebreak[4]
% p131-b21; source PDF page 131
\blockanchor{p131-b21}
\displayformula{\int_ {0} ^ {+ \infty} \frac {f (b x) - f (a x)}{x} \mathrm{d} x = f (0) \ln \frac {a}{b}.}
